\documentclass[11pt]{article}

\usepackage{../ppackage}
\usepackage{bbm}
\usepackage{import}
\usepackage{booktabs}





\usepackage{tikz}
\usetikzlibrary{arrows.meta,patterns}
\usetikzlibrary{ipe} % ipe compatibility library

\usepackage{../../../Notes/tikzit}
\usepackage{../../../Notes/utf8math}

\input{../../../Notes/TIKZ/digraph.tikzstyles}



\usepackage{url}


\newcommand{\solution}{
\bigskip\noindent
	\textbf{Solution: \\}}
	


\DeclareMathOperator{\size}{size}
\DeclareMathOperator{\conv}{conv}
\newcommand{\SV}{\mathrm{SV}}
\newcommand{\bigO}{O}
\newcommand{\cut}{\mathrm{cut}}
\newcommand{\LLL}{\mathrm{LLL}}
\newcommand{\setR}{\mathbb{R}}
\newcommand{\setZ}{\mathbb{Z}}
\newcommand{\setQ}{\mathbb{Q}}
\newcommand{\setC}{\mathbb{C}}
\newcommand{\setN}{\mathbb{N}}
\newcommand{\wt}[1]{\widetilde{#1}}
\newcommand{\opt}{{\sc 0/1-opt}\xspace}
\newcommand{\aug}{{\sc 0/1-aug}\xspace}
\newcommand{\psep}{{\sc 0/1-psep}\xspace}
\newcommand{\sep}{{\sc 0/1-sep}\xspace}
\newcommand{\fopt}{{\sc 0/1-testopt\xspace} }

\newcommand{\hpp}{\mathrm{HPP}}
\newcommand{\nodes}{\mathcal{V}}
\newcommand{\vol}{\mathrm{vol}}
\newcommand{\diag}{\mathrm{diag}}
\newcommand{\arcs}{\mathcal{A}}
\newcommand{\edges}{\mathcal{E}}
\newcommand{\paths}{\mathscr{P}}
\newcommand{\cycles}{\mathcal{C}}




\newcommand{\K}{{\mathcal K}}
\newcommand{\A}{{A}}
\newcommand{\B}{{B}}
\newcommand{\T}{\mathscr{T}}
\newcommand{\eE}{\mathscr{E}}
\newcommand{\eS}{\mathscr{S}}
\newcommand{\eP}{\mathscr{P}}
\newcommand{\eM}{\mathscr{M}}



\newcommand{\transp}{^{\mathrm{T}}}

\newcommand{\smallmat}[1]{\left( \begin{smallmatrix} #1 \end{smallmatrix}\right)}

\newcommand{\mat}[1]{ \begin{pmatrix} #1 \end{pmatrix}}
\newcommand{\smat}[1]{ \big(\begin{smallmatrix} #1 \end{smallmatrix}\big)}

\newcommand{\pc}{\mathscr{P}}
\newcommand{\ob}{\mathscr{O}}
\newcommand{\odds}{\mathscr{W}}
\newcommand{\up}{\mathscr{U}}
\newcommand{\ef}{\mathscr{F}}
\newcommand{\eh}{\mathscr{H}}
\newcommand{\ev}{\mathscr{V}}
\newcommand{\ec}{\mathscr{C}}
\newcommand{\eu}{\mathscr{U}}

\newcommand{\lex}{\mathrm{lex}}

\renewcommand{\leq}{\leqslant}
\renewcommand{\geq}{\geqslant}









\newcommand{\linhull}{\mathrm{lin.hull}}
\newcommand{\affhull}{\mathrm{affine.hull}}
\newcommand{\charcone}{\mathrm{char.cone}}
\newcommand{\cone}{\mathrm{cone}}
\newcommand{\rank}{\mathrm{rank}}
\newcommand{\wb}[1]{\overline{#1}}



\usepackage{enumerate}

      
\institute{\'Ecole Polytechnique F\'ed\'erale de Lausanne}
\lecture{Discrete Optimization}
\faculty{Prof. Friedrich Eisenbrand}
\term{Spring 2026}
\publishdate{May 12, 2026}
\duedate{ }
\problemset{Problem Set~11 (Solutions)}

\begin{document}
\makeheader

\begin{enumerate}[1)]

\item 
Let $M\in \setZ^{n\times m}$ be totally unimodular. Prove that the following matrices
are totally unimodular as well:
\begin{enumerate}[i)]
\item $M^T$
\item  $( M \quad I_n )$
\item $(M \quad -M)$
\item $M \cdot (I_n - 2 e_j e_j^T )$ for some $j$
\end{enumerate}
Here $I_n$ is the $n\times n$ identity matrix; $e_j$ is the vector having a $1$ in the $j^{th}$ component, and $0$ in the other components.

\begin{solution}
\begin{enumerate}[i)]
\item Let $A$ be a square submatrix of $M^T$ . Then
$\det(A)= \det(A^T ) ∈ \{−1,0,1\}$
as $A^T$ is a square submatrix of $M$ and $M$ is totally unimodular.

\item Let $A$ be a square submatrix of $(M \ I_n)$. Let $a_1,..., a_k$ be the columns of $A$ that originate from $I_n$ . Hence, each of these columns has as most one 1-entry, the other entries are 0. Hence, using Laplace-expansion successively along these columns we get
that $|\det(A)| = | \det(A')|$ for some square submatrix of $M$ . Since $M$ is TU, this shows
$\det(A) ∈ \{−1,0,1\}$.

\item Let $A$ be a square submatrix of $\begin{pmatrix} M&  -M \end{pmatrix}$. Let $a_1,..., a_k$ be the columns of $A$ that
originate from $−M$. Let $A'$ be the matrix obtained from $A$ by multiplying $a_1,..., a_k$ with
$-1$. Hence
$| \det(A)| = | det(A')|$.
Now we distinguish two cases.
Case 1: $A'$ is (up to permutation of columns) a square submatrix of $M$ . Since $M$ is TU,
we have $\det(A') ∈ \{−1,0,1\}$.
Case 2: $A'$ has at least two identical columns. Hence $\det(A')= 0$.
We conclude that in both cases we have $\det(A) ∈ \{−1,0,1\}$.
\item Observe that $M\cdot (I_n− 2e_j e_j^T)$ is obtained from $M$ by multiplying one column with $-1$.
Thus, $M\cdot (I_n− 2e_j e_j^T)$ is (up to permutation of columns) a submatrix of  $\begin{pmatrix}M & -M\end{pmatrix}$. Thus this matrix is also TU. 
\end{enumerate}
\end{solution}


\item A family $\mathcal{F}$ of subsets of a finite groundset $E$ is
\emph{laminar}, if for all  $C,D\in \mathcal{F}$, one of the following holds:
$$ (i)~C\cap D = \emptyset,~~(ii)~C \subseteq D,~~(iii)~D\subseteq C. $$

Let $\mathcal{F}_1$ and $\mathcal{F}_2$ be two laminar families of the same groundset $E$ and consider its union $\mathcal{F}_1\cup\mathcal{F}_2$.
Define the $|\mathcal{F}_1\cup\mathcal{F}_2|\times |E|$ adjacency matrix $A$ as follows:
For $F\in\mathcal{F}_1\cup\mathcal{F}_2$ and $e\in E$ we have
$A_{F,e}=1$, if $e\in F$ and $A_{F,e}=0$ otherwise.

Show that $A$ is totally unimodular.

\begin{solution}
Let $F_1$ and $F_2$ be two laminar families on the same groundset $E$ , and $A$ the corresponding
adjacency matrix. Observe that every square submatrix of $A$ also is the adjacency matrix of
two laminar families: Removing a row from $A$ corresponds to deleting a set from the laminar
families. Removing a column from $A$ corresponds to removing an element of the ground-
set from all sets of the laminar families. Both operations preserve the structure of laminar
families.


For that reason, it is sufficient to show the following statement: Every square matrix $A$ that
is adjacency matrix of two laminar families has determinant $±1$ or $0$.
We will transform $A$ with elementary column operations as follows: Let $e ∈ E$ be an element from the groundset that is contained in at least two sets from $F_1$. Let $S_1,...,S_k$ be
the sets of $F_1$ with $e ∈ S_i$ . Using the properties of laminar families, we know that there is a
$l ∈ \{1,...,k\}$ such that $S_l ⊆ S_i$ for all $i= 1,...,k$. Redefine $S_i := S_i \setminus S_l$ for all $i\neq l$ . Observe
that $F_1$ is still a laminar family after this modification. Also observe that the operation of removing the set $S_l$ corresponds to subtraction the row $S_l$ from the other rows $S_i$ in the matrix
A.
Hence we can apply this transformation until each $e ∈ E$ is contained in at most one set of
$F_1$. Similarly we can apply this transformation to $F_2$ until each $e ∈ E$ is contained in at most
one set of $F_2$. Applying the corresponding elementary row operations to A yields a matrix
$A'$ with $\det(A)= \det(A')$.  $A'$ has the property that there are two disjoint subsets of rows, the
rows corresponding to $F_1$ and the rows corresponding to $F_2$, such that each column of $A'$
has at most one 1-entry in the rows of $F_1$ and at most one 1-entry in the rows of $F_2$. All other
entries are $0$.\\
Let $A''$ be the submatrix of $A'$ consisting only of the columns with two 1-entries. Note
that this is a node-edge incidence matrix of a bipartite graph. Hence $A''$ is TU. With Exercise 1.2 we get that $A'$ is TU. Hence $\det(A) ∈ \{−1,0,1\}$.
\end{solution}

\item Let $G$ be a graph and let $A$ be its node-edge incidence matrix. We have seen that if $G$ is
bipartite then $A$ is totally unimodular. Prove the converse, i.e., if $A$ is totally unimodular then $G$
is bipartite.

\begin{solution}
Let the incidence matrix of $G$ be totally unimodular and assume towards contradiction that $G$ is not bipartite. Then $G$ must contain a cycle of odd length. Let this cycle contain some vertices of $G$ $\{v_1, \hdots, v_{2k+1}\}$ for some $k \in \mathbb{N}$. Let the edges of this cycle be $\{e_1, \hdots, e_{2k+1}\}$. Now, consider the submatrix of $A$ indexed by $[v_1, \hdots, v_{2k+1}] \times [e_1, \hdots, e_{2k+1}]$.  
Then this submatrix of the cycle (up to permutation of the columns) looks as follows: \\

\begin{tabular}{l | l l l l l }
 & $e_1$ & $e_2$ &$e_3$ & $\hdots$ & $e_{2k+1}$ \\
 \midrule
$v_1$ & 1 & 0 & 0 &  $\hdots$ & 0 \ 1 \\
$v_2$ & 1 & 1 & 0 & $\hdots$ & 0 \ 0 \\
$v_3$ & 0 & 1 & 1 & $\hdots$ &  0 \ 0  \\
$\vdots$ & & $\ddots$ & & &  \\
$v_{2k+1}$ & 0 & 0 & 0 & $\hdots$ & 1 \ 1
\end{tabular}

Then since the number of rows and columns is odd, we can do row reduction and end with one row that has a value 2 giving the whole submatrix a determinant of 2. This means $A$ is not submodular, a contradiction. 
\end{solution}

% \item Which of these two matrices is totally unimodular? Justify your answer. 

% \begin{align*}
%   i).&
%     \begin{pmatrix}
%       1 & 1  & 0 &  1 & 0 \\
%       0 & 0 &  1 &  0 &  1 \\
%       0 & 1 &  1 &  1 & 0 \\
%       1 & 1 &  0 &  0 & 0 \\
%       0 & 1 &  0 &  0 & 1
%     \end{pmatrix}
%     \\
%     \vspace{3mm}
%     \\
%   ii).&
%     \begin{pmatrix}
%       1&  1& 1& 1& 1 \\
%       1& 1& 0& 0& 0 \\
%       1& 1& 1& 0& 0 \\
%       1&  1& 1& 1& 0
%     \end{pmatrix}
% \end{align*}

\item Sudoku is the following puzzle: Given a matrix
  \begin{displaymath}
    A ∈ \{1,\dots,9,X\}^{9 ×9} 
  \end{displaymath}
  the task is to replace each $X$ in $A$ by a number in $\{1,\dots,9\}$ such that the following holds.
  \begin{enumerate}[a)] 
  \item Each line of $A$ contains all numbers in  $\{1,\dots,9\}$
  \item Each column of $A$ contains all numbers in  $\{1,\dots,9\}$
  \item If $A$ is written as
    \begin{displaymath}
      A =
      \begin{pmatrix}
        B_{11} & B_{12} & B_{13} \\
        B_{21} & B_{22} & B_{23} \\
        B_{31} & B_{32} & B_{33} 
      \end{pmatrix}
    \end{displaymath}
    where the $B_{ij}$ are $3 ×3$ matrices, the each of them contains all numbers in  $\{1,\dots,9\}$.

    \bigskip \noindent

    In the following, you are asked to design an integer program whose feasible solutions are solutions of a given Sudoku. The integer program has $9^3$ variables
    \begin{displaymath}
      x_{i,j,k} =
      \begin{cases}
        1 & \text{ if } k \text{ is the number in cell } (i,j),  \\
        0 & \text{ otherwise.} 
      \end{cases}
    \end{displaymath}


    The following set of constraints guarantees that each cell of the solution contains a number
    \begin{displaymath}
      1≤i,j≤9: \quad  ∑_{k=1}^9 x_{i,j,,k} = 1. 
    \end{displaymath}

    Write now constraints that guarantee the following.
    \begin{itemize}
    \item Each row must contain each number.
    \item Each column must contain each number.
    \item Each $B_{ij}$ must contain each number. 
    \end{itemize}

    Describe the final integer programming problem that links some of the variables to the input Sudoku. 
    
  \end{enumerate}

\begin{solution}
The condition that each row must contain each number, corresponds to the set of constraints
\[
\sum_{j=1}^9 x_{i,j,k} = 1,\quad \forall\, 1 \leq i,k \leq 9.
\]
The condition that each column must contain each number, corresponds to the set of constraints
\[
\sum_{i=1}^9 x_{i,j,k} = 1,\quad \forall\, 1 \leq j,k \leq 9.
\]
The condition that each $B_{ij}$ must contain each number, corresponds to the set of constraints
\[
\sum_{\alpha=0}^2 \sum_{\beta=0}^2 x_{i+\alpha,j+\beta,k} = 1,\quad \forall\, i,j=1,4,7,\; 1\leq k\leq 9.
\]
To get the final integer program, we only need the set of constraints which guarantees that the solution is compatible with the input, i.e., 
\[
x_{i,j,k} = 1,\quad \text{if in the input of Sudoku, the cell}\; (i,j) \;\text{is assigned the number} \; k
\]
and the integrality constraints
\[
x_{i,j,k} \in \{0,1\},\quad \forall\, 1 \leq i,j,k \leq 9.
\]
\end{solution}

\item \label{e:weighted-vextex-cover}
\emph{Bipartite vertex cover with demands} is the following problem. We are given a bipartite graph $G = (A\cup B,E)$ and edge costs $w: E ⟶ ℕ_0$. The \emph{weight} of a matching  $M ⊆ E$ is defined as  
  \begin{displaymath}
   w(M) = ∑_{e∈M} w_e. 
  \end{displaymath}
  A \emph{vertex cover for demands} $w$  for this setting is an integer vector $x ∈ ℕ_0^{|V|}$ such that
    \begin{displaymath}
      x_u + x_v ≥ w_{uv}, \text{ for each edge } uv ∈ E. 
    \end{displaymath}
    Show the following version of König's theorem:
    \begin{quote}
      The maximum weight of a matching of $G$ is equal to the minimum $ℓ_1$-norm of a weighted vertex cover. 
    \end{quote}

\begin{solution}
The linear program for the maximum weighted matching problem is
\begin{align*}
\max \quad & \sum_{e \in E} w_e y_e \\
\text{s.t.} \quad & \sum_{e \in \delta(v)} y_e \leq 1, \quad \forall v \in A \cup B \\
& y_e \geq 0, \quad \forall e \in E
\end{align*}
where $\delta(v)$ is the set of edges incident to $v$. 
Its dual linear program encodes the minimum vertex cover for demands problem and is given by
\begin{align*}
\min \quad & \sum_{v \in A \cup B} x_v \\
\text{s.t.} \quad & x_u + x_v \geq w_{uv}, \quad \forall uv \in E \\
& x_v \geq 0, \quad \forall v \in A \cup B
\end{align*}
As shown in the lecture, the constraint matrix of the primal LP is totally unimodular. Then both the primal and the dual LP have integral optimal solutions (note that $w$ is integral). The statement follows from strong duality.

\end{solution}

\end{enumerate}

  

\end{document}

%%% Local Variables:
%%% mode: latex
%%% TeX-master: t
%%% End:
